The Process class is the primary handle for operating on a single target process.
Process objects may be created by calls to the static functions
Process::createProcess or Process::attachProcess, or in response to
certain types of events (e.g, fork on UNIX systems).

The static functions of the Process class serve as a central location for
performing general ProcControlAPI operations, such as handleEvents and
registerEventCallback when dealing with callbacks. 

\begin{apient}
Process::ptr Process::const_ptr
\end{apient}

\apidesc{
The Process::ptr and Process::const\_ptr respectively represent a pointer and a
const pointer to a Process object. Both pointer types are reference counted
and will cause the underlying Process object to be cleaned when there are no
more references. ProcControlAPI will maintain internal references to any
Process it actively controls, relinquishing those references when the process
either exits or is detached.
}

\begin{apient}
enum Process::cb_action_t {
  cbDefault,
  cbThreadContinue,
  cbThreadStop, 
  cbProcContinue,
  cbProcStop
}
\end{apient}

\apidesc{
}

\begin{apient}
struct Process::cb_ret_t {
  cb_ret_t(cb_action_t p)
  cb_ret_t(cb_action_t p, cb_action_t c)
\end{apient}

\apidesc{
The cb\_ret\_t enum is used as the return type for callback functions registered
through Process::registerEventCallback(). A callback function can specify
whether the thread or process associated with its event should be stopped or
continued by respectively returning cbThreadContinue, cbThreadStop,
cbProcContinue, or cbProcStop. The cbDefault return value returns a Process
and Thread to the original state before the event occurred.
Some events, such as process spawn or thread create involve two processes or
threads. In this case the ProcControlAPI user can specify a cb\_action\_t
value for both the parent and child using the two parameter constructor for
cb\_ret\_t.
}

\begin{apient}
typedef Process::cb_ret_t(*cb_func_t)(Event::const_ptr)
\end{apient}

\apidesc{
The cb\_func\_t type is a function pointer type for functions that can handle
event callbacks. The callback function gets an Event::const\_ptr as input,
which points to the Event that triggered the callback. The cb\_func\_t
function should return a cb\_ret\_t describing what to do with the process
after handling the event.
}

\begin{apient}
enum OSType{
  OSNone,
  Linux,
  FreeBSD,
  Windows
}
\end{apient}

\apidesc{
A value from this enum is returned from Process::getOS and signifies the current
OS on which the target process is running.
}

\begin{apient}
typedef std::pair<Dyninst::Address, Dyninst::Address> MemoryRegion
\end{apient}

\apidesc{
The MemoryRegion type represents a region of allocated memory, and the first
part of the pair is the start address, the second, the end.
}

\begin{apient}
static Process::ptr createProcess(std::string executable,
                                  const std::vector<std::string> &argv,
                                  const std::vector<std::string> &envp = emptyEnv,
                                  const std::map<int,int> &fds = emptyFDs)
\end{apient}

\apidesc{
This function creates a new process by launching an executable file named by
executable with the arguments specified by argv, the environment specified in
envp, and it returns a pointer to the new Process object upon success. The
new process will be created with its initial thread in the stopped state. 
It is an error to call this function from a callback.
If the fds map is not empty, then the new process will be created with the file
descriptors from the fds\'  first elements dup2 mapped to
the file descriptors in fds\'  second elements.
If envp is empty, the environment will be inherited from the calling process.
ProcControlAPI may deliver callbacks when this function is called.
This function returns Process::ptr() on error, and a subsequent call to
getLastError returns details on the error.
}

\begin{apient}
static Process::ptr attachProcess(Dyninst::PID pid,  std::string executable = "")
\end{apient}

\apidesc{
This function creates a new Process object by attaching to the PID specified by
pid. The new Process object will be returned from this function upon success.
The executable argument is optional, and can be used to assist ProcControlAPI
in finding the process\'  executable on operating systems
where this cannot be easily determined (currently on AIX). The new process
will be returned with all of its threads in the stopped state.
It is an error to call this function from a callback.
ProcControlAPI may deliver callbacks when this function is called.
This function return Process::ptr() on error, and a subsequent call to
getLastError returns details on the error.
}

\begin{apient}
static bool handleEvents(bool block)
\end{apient}

\apidesc{
This function causes ProcControlAPI to handle any pending debug events and
deliver callbacks. When an event requires a
callback ProcControlAPI needs control of the main thread in order to deliver
the callback. This function gives control of the main thread to
ProcControlAPI for callback delivery. A user can know when to call
handleEvents by using the EventNotify interface. 
If the block parameter is true, then handleEvents will block until at least one
debug event has been handled. If block is false then handleEvents returns
immediately if no events are ready to be handled. 
This function returns true if it handled at least one event and false otherwise.
It is an error to call this function from a callback.
}

\begin{apient}
static bool registerEventCallback(EventType evt, cb_func_t cbfunc)
\end{apient}

\apidesc{
This function registers a new callback function with ProcControlAPI. Upon
receiving an event with type evt, ProcControlAPI will deliver a callback with
that event to the cbfunc function. Multiple functions can be registered to
receive callbacks for a single EventType, and a single function can be
registered with multiple EventTypes. 
If multiple callback functions are registered with a single EventType, then it
is undefined what order those callback functions will be invoked in. In this
case the cb\_ret\_t result of the last callback function called will be used to
determine what stop or continue operations should be performed on the process.
If a single callback function is registered for the same EventType multiple
times, then ProcControlAPI will only invoke one call to the callback function
for each instance of the EventType.
This function returns true on success and false on error. Upon an error a
subsequent call to getLastError returns details on the error.
}

\begin{apient}
static bool removeEventCallback(EventType evt, cb_func_t cbfunc)
\end{apient}

\apidesc{
This function un-registers a callback that was registered with
registerEventCallback. After a successful call to this function the callback
function cbfunc will stop being called for events with EventType evt. Other
callback functions registered for evt will not be affected. Other instances
of cbfunc registered for different EventTypes will not be affected.
This function returns true if a callback was successfully removed and false
otherwise. Upon an error a subsequent call to getLastError returns details on
the error.
}

\begin{apient}
static bool removeEventCallback(EventType evt)
\end{apient}

\apidesc{
This function unregisters all callback functions associated with the EventType
evt. After a successful call to this function ProcControlAPI will stop
delivering callbacks for evt until a new callback function is registered.
This function returns true if a callback was successfully removed and false
otherwise. Upon an error a subsequent call to getLastError returns details on
the error.
}

\begin{apient}
static bool removeEventCallback(cb_func_t func)
\end{apient}

\apidesc{
This function unregisters all instances of callback function func from any
callback with any EventType. 
This function returns true if a callback was successfully removed and false
otherwise. Upon an error a subsequent call to getLastError returns details on
the error.
}

\begin{apient}
Dyninst::PID getPid() const
\end{apient}

\apidesc{
This function returns an OS handle referencing the process. On UNIX systems
this is the pid of the process. 
}

\begin{apient}
Dyninst::Architecture getArchitecture() const
\end{apient}

\apidesc{
This function returns an enum that describes the architecture of the target
process.
}

\begin{apient}
Dyninst::OSType getOS () const
\end{apient}

\apidesc{
This function returns an enum that describes the OS of the target process. See
the beginning of this section for the definition of Dyninst::OSType.
}

\begin{apient}
bool supportsLWPEvents () const
\end{apient}

\apidesc{
This function returns true if the target process can throw LWP create and
destroy events and false otherwise.
}

\begin{apient}
bool supportsUserThreadEvents() const
\end{apient}

\apidesc{
This function returns true if the target process can throw user thread create
and destroy events and false otherwise.
}

\begin{apient}
bool supportsFork () const
\end{apient}

\apidesc{
This function returns true if the fork system call is supported in the target
process and false otherwise.
}

\begin{apient}
bool supportsExec () const
\end{apient}

\apidesc{
This function returns true if the exec system call is supported in the target
process and false otherwise.
}

\begin{apient}
bool isTerminated() const
\end{apient}

\apidesc{
This function returns true if the target process has terminated (either via a
crash or normal exit) or if the ProcControlAPI has detached from the target
process. It returns false otherwise.
}

\begin{apient}
bool isExited() const
\end{apient}

\apidesc{
This function returns true of the target process exited via a normal exit (e.g,
calling the exit function or returning from main). It returns false
otherwise.
}

\begin{apient}
int getExitCode() const
\end{apient}

\apidesc{
If a target process exited normally then this function returns its exit code.
The return result of this function is undefined if the
Process\' isExited function returns false.
}

\begin{apient}
bool isCrashed() const
\end{apient}

\apidesc{
This function returns true if the target process exited because of a crash. It
returns false otherwise.
}

\begin{apient}
int getCrashSignal() const
\end{apient}

\apidesc{
If a target process exited because of a crash, then this function returns the
signal that caused the target process to crash. The return result of this
function is undefined if the Process\'  isCrashed function
returns false.
}

\begin{apient}
bool hasStoppedThread() const
\end{apient}

\apidesc{
This function returns true if the target process has at least one thread in the
stopped state. It returns false otherwise or if an error occurs. In the
event of an error a call to getLastError returns details on the error.
}

\begin{apient}
bool hasRunningThread() const
\end{apient}

\apidesc{
This function returns true if the target process has at least one thread in the
running state. It returns false otherwise or if an error occurs. In the
event of an error a call to getLastError returns details on the error.
}

\begin{apient}
bool allThreadsStopped() const
\end{apient}

\apidesc{
This function returns true if all threads in the target process are in the
stopped state. It returns false otherwise or if an error occurs. In the
event of an error a call to getLastError returns details on the error.
}

\begin{apient}
bool allThreadsRunning() const
\end{apient}

\apidesc{
This function returns true if all threads in the target process are in the
running state. It returns false otherwise or if an error occurs. In the
event of an error a call to getLastError returns details on the error.
}

\begin{apient}
bool allThreadsRunningWhenAttached() const
\end{apient}

\apidesc{
This function returns true if all threads were running when the controller
process attached to this process. It returns false if any threads were stopped.
If the target process was created instead of attached, this function returns
true.
}

\begin{apient}
bool continueProc()
\end{apient}

\apidesc{
This function will move all threads in the target process into the running
state. This function returns true if at least one thread was continued as
part of the call, and false otherwise.
It is an error to call this function from a callback.
ProcControlAPI may deliver callbacks when this function is called.
This function return false on error, and a subsequent call to getLastError
returns details on the error.
}

\begin{apient}
bool stopProc()
\end{apient}

\apidesc{
This function will move all threads in the target process into the stopped
state. This function returns true if at least one thread was stopped as part
of the call, and false otherwise.
It is an error to call this function from a callback.
ProcControlAPI may deliver callbacks when this function is called.
This function return false on error, and a subsequent call to getLastError
returns details on the error.
}

\begin{apient}
bool detach(bool leaveStopped = false)
\end{apient}

\apidesc{
This function will detach ProcControlAPI from the target process.
ProcControlAPI will no longer be able to control or receive events from the
target process. All breakpoints will be removed from the target. This
function returns true on success and false on error. Upon an error a
subsequent call to getLastError returns details on the error.
If the leaveStopped parameter is set to true, and the process is in a stopped
state, then the target process will be left in a stopped state after the
detach. 
It is an error to call this function from a callback.
}

\begin{apient}
bool temporaryDetach()
\end{apient}

\apidesc{
This function temporarily detaches from the target process, but leaves the
Process data structure intact. This functionality is commonly called
detach-on-the-fly. The target process will not report new events nor be
controllable or able to be queried by the user. Breakpoints are removed from
the process. The reAttach function will reconnect the process after this
call.
This function returns true on success and false upon error.
It is an error to call this function from a callback.
}

\begin{apient}
bool reAttach()
\end{apient}

\apidesc{
This function reconnects to the target process after a temporaryDetach call.
Any breakpoints will be re-inserted back into the function, and if threads
have been created or destroyed during the time detached new events will be
thrown for them.
This function returns true on success and false upon error.
It is an error to call this function from a callback.
}

\begin{apient}
bool terminate()
\end{apient}

\apidesc{
This function forcefully terminated the target process. Upon a successful call
to this function the target process will end execution. The Process object
will record the target process as having crashed. This function returns true
on success and false on error. Upon an error a subsequent call to
getLastError returns details on the error.
It is an error to call this function from a callback.
}

\begin{apient}
const ThreadPool &threads()
const ThreadPool &threads()
\end{apient}

\apidesc{
Returns the \code{ThreadPool} of the process.
}

\begin{apient}
const LibraryPool &libraries() const
LibraryPool &libraries()
\end{apient}

\apidesc{
Returns the \code{LibraryPool} of the process.
}

\begin{apient}
bool addLibrary(std::string libname)
\end{apient}

\apidesc{
This function causes the specified library to be loaded into the process. It
will trigger an event (and thus a user callback) for each library loaded
(including dependencies). 
}

\begin{apient}
void *getData () const
void setData (void *p)
\end{apient}

\apidesc{
These functions respectively get and set an opaque data object that can be
associated with this process. The data is not interpreted by ProcControlAPI,
but remains associated with the process.
}

\begin{apient}
unsigned getMemoryPageSize() const
\end{apient}

\apidesc{
This function returns memory page size for the current OS on which the target
process is running.
}

\begin{apient}
Dyninst::Address mallocMemory(size_t size)
Dyninst::Address mallocMemory(size_t size, Dyninst::Address addr)
\end{apient}

\apidesc{
These functions allocate a region of memory in the target
process\'  address space of size size. Upon a successful
call these functions will map an area of memory in the target process that is
readable, writeable and executable. The \code{mallocMemory(size\_t)} function will
allocate memory at any available address. The \code{mallocMemory(size\_t, Dyninst::Address)}
function will only allocate memory at the specified address, addr.
It is an error to call this function from a callback.
ProcControlAPI may deliver callbacks when this function is called.
Upon success these functions return the start address of memory that was
allocated and 0 otherwise. Upon an error a subsequent call to getLastError
returns details on the error.
}

\begin{apient}
bool freeMemory(Dyninst::Address addr)
\end{apient}

\apidesc{
This function will free a region of memory that was allocated by the
mallocMemory function. Upon a successful call to this function, the area of
memory starting at addr will be unmapped and no longer accessible to the target
process. It is an error to call this function with an address that was not
returned by mallocMemory.
It is an error to call this function from a callback.
ProcControlAPI may deliver callbacks when this function is called.
Upon success this function returns true, otherwise it returns false. Upon an
error a subsequent call to getLastError returns details on the error.
}

\begin{apient}
bool writeMemory(Dyninst::Address addr, void *buffer, size_t size) const
\end{apient}

\apidesc{
This function writes to the target process\' s memory. The
addr parameter specifies an address in the target process to which
ProcControlAPI should write. The buffer and size parameters specify a region
of controller process memory that will be copied into the target process.
It is an error to call this function on a Process that does not have at least
one Thread in a stopped state.
This function returns true on success and false on error. Upon an error a
subsequent call to getLastError returns details on the error.
}

\begin{apient}
bool readMemory(void *buffer, Dyninst::Address addr, size_t size) const
\end{apient}

\apidesc{
This function reads from the target process\'  memory. The
addr and size parameters specify an address in the target process from which
ProcControlAPI should read. The buffer parameter specifies an address in the
controller process where ProcControlAPI should write the copied bytes. 
It is an error to call this function on a Process that does not have at least
one Thread in a stopped state.
This function returns true on success and false on error. Upon an error a
subsequent call to getLastError returns details on the error.
}

\begin{apient}
bool getMemoryAccessRights(Dyninst::Address addr, mem_perm& rights)

bool setMemoryAccessRights(Dyninst::Address addr, size_t size, mem_perm rights, mem_perm& oldRights)
\end{apient}

\apidesc{
These functions respectively get and set memory permission at the specified
address, addr. The setMemoryAccessRights function also affects a region of
memory in the target process\' s address space of size size.
}

\begin{apient}
bool findAllocatedRegionAround(Dyninst::Address addr, MemoryRegion& memRegion)
\end{apient}

\apidesc{
This function finds a region of allocated memory memRegion contains address
addr, and returns true on success, otherwise false.
}

\begin{apient}
bool addBreakpoint(Dyninst::Address addr, Breakpoint::ptr bp) const
\end{apient}

\apidesc{
This function inserts the Breakpoint specified by bp into the target process at
address addr.
It is an error to call this function on a Process that does not have at least
one Thread in a stopped state.
This function returns true on success and false on error. Upon an error a
subsequent call to getLastError returns details on the error.
}

\begin{apient}
bool rmBreakpoint(Dyninst::Address addr, Breakpoint::ptr bp) const
\end{apient}

\apidesc{
This function removes the Breakpoint specified by bp at address addr from the
target process.
This function returns true on success and false on error. Upon an error a
subsequent call to getLastError returns details on the error.
}

\begin{apient}
bool postIRPC(IRPC::ptr irpc) const
\end{apient}

\apidesc{
This function posts the given irpc to the Process. ProcControlAPI selects a
Thread from the Process to run the iRPC and put irpc into that
Thread\' s queue of posted IRPCs.
Each instance of an IRPC object can be posted at most once. It is an error to
attempt to post a single IRPC object multiple times.
This function returns true on success and false on error. Upon an error a
subsequent call to getLastError returns details on the error.
}

\begin{apient}
bool runIRPCSync(IRPC::ptr irpc)
\end{apient}

\apidesc{
This function posts an irpc, similar to Process::postIRPC; continues the thread
the irpc was posted to; and returns when the irpc has completed running. The
thread will be returned to its original running state when this function
returns.
This function returns true if the irpc was successfully run, and false
otherwise. Note that stopping the thread that is running the irpc while this
function waits for irpc completion causes this function to return an error.
It is an error to call this function from a callback.
}

\begin{apient}
bool runIRPCAsync(IRPC::ptr irpc)
\end{apient}

\apidesc{
This function posts an irpc, similar to Process::postIRPC, and then continues
the thread the irpc was posted to.
This function returns true if the irpc was successfully posted and run, and
false otherwise.
It is an error to call this function from a callback.
}

\begin{apient}
bool getPostedIRPCs(std::vector<IRPC::ptr> &rpcs) const
\end{apient}

\apidesc{
This function returns all IRPCs posted to this Process in the rpcs vector.
This list does not include any IRPCs currently
running- see Thread::getRunningIRPC() for this
functionality.
This function returns true on success and false on error. Upon an error a
subsequent call to getLastError returns details on the error.
}

\begin{apient}
LibraryTracking *getLibraryTracking()
ThreadTracking *getThreadTracking()
LWPTracking *getLWPTracking()
FollowFork *getFollowFork()
SignalMask *getSignalMask()
\end{apient}

\apidesc{
These functions return pointers to configuration objects for platform-specific
features associated with this Process object. 
These functions return NULL if the specified feature is unsupported on the
current platform.
}
